% ============================================================
% Reflective Intelligent Bodies (RIB)
% Companion to q_complement.tex and Part III of monograph_v2.tex
% Author: R. Wydaeghe
% ============================================================

\documentclass[11pt,a4paper]{article}

\usepackage[margin=2.4cm]{geometry}
\usepackage[utf8]{inputenc}
\usepackage[T1]{fontenc}
\usepackage{lmodern}
\usepackage{microtype}
\linespread{1.04}

\usepackage{amsmath,amssymb,amsthm,mathtools}
\usepackage{bm}
\usepackage{mathrsfs}
\usepackage{bbm}
\usepackage{cancel}
\usepackage{xfrac}
\allowdisplaybreaks[2]

\usepackage{empheq}
\usepackage[most]{tcolorbox}
\tcbset{highlight math style={%
    enhanced, colframe=black!70, colback=gray!6,
    arc=2pt, boxrule=0.5pt, left=3pt, right=3pt}}

\usepackage{siunitx}
\sisetup{detect-all, range-phrase = {--}, range-units = single,
         per-mode = symbol, inter-unit-product = {\,}}

\usepackage{enumitem}
\usepackage{xcolor}
\usepackage{xspace}
\usepackage{fancyhdr}
\pagestyle{fancy}
\fancyhf{}
\cfoot{\thepage}
\renewcommand{\headrulewidth}{0pt}

\usepackage{titlesec}
\titleformat{\section}{\large\bfseries}{\thesection}{1em}{}
\titleformat{\subsection}{\normalsize\bfseries}{\thesubsection}{1em}{}

\usepackage[colorlinks=true, linkcolor={blue!55!black},
            citecolor={green!45!black}, urlcolor={blue!65!black}]{hyperref}
\usepackage[nameinlink,capitalize]{cleveref}

\usepackage{aliascnt}
\theoremstyle{definition}
\newtheorem{theorem}{Theorem}[section]
\newaliascnt{proposition}{theorem}
\newtheorem{proposition}[proposition]{Proposition}
\aliascntresetthe{proposition}
\newaliascnt{lemma}{theorem}
\newtheorem{lemma}[lemma]{Lemma}
\aliascntresetthe{lemma}
\newaliascnt{corollary}{theorem}
\newtheorem{corollary}[corollary]{Corollary}
\aliascntresetthe{corollary}
\newaliascnt{definition}{theorem}
\newtheorem{definition}[definition]{Definition}
\aliascntresetthe{definition}
\newaliascnt{remark}{theorem}
\newtheorem{remark}[remark]{Remark}
\aliascntresetthe{remark}
\crefname{proposition}{Proposition}{Propositions}
\Crefname{proposition}{Proposition}{Propositions}
\crefname{lemma}{Lemma}{Lemmas}\Crefname{lemma}{Lemma}{Lemmas}
\crefname{corollary}{Corollary}{Corollaries}\Crefname{corollary}{Corollary}{Corollaries}
\crefname{definition}{Definition}{Definitions}\Crefname{definition}{Definition}{Definitions}
\crefname{remark}{Remark}{Remarks}\Crefname{remark}{Remark}{Remarks}

% Macros aligned with monograph and q_complement
\newcommand{\khat}{\hat{\bm{k}}}
\newcommand{\nhat}{\hat{\bm{n}}}
\newcommand{\rhat}{\hat{\bm{r}}}
\newcommand{\rr}{\mathbf{r}}
\newcommand{\EE}{\mathbf{E}}
\newcommand{\Sinc}{S_{\mathrm{inc}}}
\newcommand{\Sab}{S_{\mathrm{ab}}}
\newcommand{\Sref}{S_{\mathrm{re}}}
\newcommand{\Sinn}{S_{\mathrm{in}}}
\newcommand{\Smis}{S_{\mathrm{mi}}}
\DeclareMathOperator{\ReLU}{ReLU}
\DeclareMathOperator{\RE}{Re}
\DeclareMathOperator{\IM}{Im}
\DeclareMathOperator{\tr}{tr}
\newcommand{\diff}{\mathrm{d}}
\newcommand{\ident}{\mathbbm{1}}
\newcommand{\Reals}{\mathbb{R}}
\newcommand{\Complex}{\mathbb{C}}
\newcommand{\half}{\tfrac{1}{2}}
\newcommand{\Gtil}{\tilde{\mathbf{G}}}
\newcommand{\gtil}{\tilde{\mathbf{g}}}
\newcommand{\Fn}{\mathbf{F}_{\!n}}
\newcommand{\Rn}{\mathbf{R}_{\!n}}
\newcommand{\Qabs}{\mathbf{Q}_{\mathrm{ab}}}
\newcommand{\Qref}{\mathbf{Q}_{\mathrm{re}}}
\newcommand{\Qinn}{\mathbf{Q}_{\mathrm{in}}}
\newcommand{\Qmis}{\mathbf{Q}_{\mathrm{mi}}}

% RIB-specific macros (subscripts, not superscripts, to avoid `\kin^*` collisions)
\newcommand{\kin}{\hat{\bm{k}}_{\mathrm{i}}}
\newcommand{\kout}{\hat{\bm{k}}_{\mathrm{o}}}
% Body scattering matrix; per-path index uses superscript-(n) to keep
% one atom and avoid double-subscript collisions.
\newcommand{\Frib}{\mathbf{F}_{\mathrm{b}}}
\newcommand{\sigb}{\sigma_{\mathrm{b}}}
\newcommand{\Qrib}{\mathbf{Q}_{\mathrm{re}}}
\newcommand{\Qbist}{\mathbf{Q}_{\mathrm{bi}}}
% JSAC macros pulled in for Sec.~\ref{sec:secondbounce}
\newcommand{\hh}{\mathbf{h}}
\newcommand{\Jmat}{\mathbf{J}}
\newcommand{\bvec}{\mathbf{b}}

\title{\textbf{Reflective Intelligent Bodies:\\
the human body as an analytic far-field scatterer}}
\author{Robin Wydaeghe\\
 Ghent University \& imec\\
 \texttt{robin.wydaeghe@ugent.be}}
\date{\today}

\makeatletter
\renewcommand{\maketitle}{%
  \thispagestyle{plain}%
  \begin{center}
    \rule{\linewidth}{1pt}\\[0.4em]
    {\LARGE\bfseries \@title}\\[0.5em]
    {\large \@author}\\[0.3em]
    {\small \@date}\\[0.4em]\rule{\linewidth}{1pt}
  \end{center}
  \vspace{0.5em}
}
\makeatother

\begin{document}
\maketitle

\begin{abstract}
A digital-twin network (DTN) that maintains a per-user body twin for
exposure compliance is, almost for free, also maintaining the
information needed to predict each body's far-field scattering
pattern.  This note develops that observation into a deployable
addition to the JSAC body twin.  Treating each body as a
Kirchhoff--physical-optics scatterer at mmWave and above, we define a
bistatic body scattering matrix $\Frib(\kin,\kout)$, a directional
exposure operator $\Qbist(\kout)$, and the body radar cross-section
$\sigb(\kin,\kout)$.  The pseudo-Brewster eigenvector locking from
\texttt{q\_complement.tex} extends to a \emph{per-direction}
statement, and that produces the centrepiece result of this note:
the capacity gain a $K$-body RIB-network can deliver is fundamentally
upper-bounded by the array's exposure budget on those bodies.
Capacity-exposure is a single Pareto frontier, not two axes.  In the
JSAC paper's plaza scenes, this yields a model-based explanation of
the empirical effective rank $3{-}4$ at $99\%$ trace.  Like the RIS
literature, RIB does not need first-principles fidelity to deploy; it
inherits the same path-amplitude calibration the JSAC twin already
performs.  What is new is the body's status in the propagation graph:
not a blocker, not a black box, but a self-tuning passive RIS whose
only knob is gait, and whose pattern the DTN already pays the
compute to predict.
\end{abstract}

\vspace{0.3em}
\noindent\textit{Style note.}~Like its sibling
\texttt{q\_complement.tex}, this is a wandering companion, not a
theorem-proof paper.  The Kirchhoff--PO derivation is rigorous in
the high-frequency regime where mmWave bodies live; everything
else is flagged as either approximate or conjectural.  The intended
home for the deployable parts of this note is a section of the
JSAC SI submission on Digital Twins for Wireless Networks; the
intended home for the loose ends is a follow-up.

\tableofcontents
\bigskip

% ============================================================
\section{Why bother, given the JSAC paper already has a body twin}
% ============================================================

The honest version of ``why this'' is short.  The JSAC paper's twin
already maintains, per served user $u$, a per-body exposure operator
$\mathbf{Q}^{(u)}$ that the precoder solver uses to enforce
compliance.  The same surface integral that produces
$\mathbf{Q}^{(u)}$ produces, with one extra direction-resolved sum,
the body's bistatic scattering matrix $\Frib^{(u)}(\kout)$.  For the
twin, RIB is a one-line addition to a code path that is already
instrumented; for the propagation graph, it lifts the body from
``blocker'' to ``analytic node'' on the same compute budget.  We do
not have to argue from scratch that body-aware ray tracing is worth
its weight in cycles; the cycles are already being spent.

Three things this buys, none of which the served-user channel
estimate $\hh$ delivers on its own.

\begin{enumerate}[leftmargin=1.6em,itemsep=3pt]
\item \textbf{A capacity-exposure Pareto bound.}  Per-direction
pseudo-Brewster locking (\cref{prop:perdir} below) ties scattered
power to absorbed power eigenvector-by-eigenvector.  Consequence:
the rank lift the body-network gives the array is
\emph{fundamentally} bounded by the array's exposure budget.  Two
axes the JSAC paper currently treats as independent are in fact one.
\Cref{thm:pareto} below makes this precise.
\item \textbf{Bystander twin without pilots.}  Reciprocity
delivers $\hh$ to served users.  It does not deliver anything for
tier-C bystanders, who are observed only by what they scatter.
RIB gives the twin a model-based predictor of that scatter, which
turns ``binary CFAR detect'' into ``predict where the body is, fit
its pose to the residual.''  This is the JSAC SI scope item
\textit{Twin-in-the-loop architectures for real-time feedback and
control} delivered for an unmonitored agent.
\item \textbf{Empirical rank explained.}  The JSAC paper
observes effective rank $3{-}4$ at $99\%$ trace in the per-user
$\mathbf{Q}^{(u)}$ at plaza scale, attributed to ``low-rank
generically''.  RIB gives a model-based reason: it is the
projection of the array's $M$-dimensional steering manifold onto
the body's angular-mode subspace, which has dimension
$M_{\mathrm{eff}}^{(u)}\sim 10{-}30$ at \SI{28}{GHz}.  An empirical
observation gets a derivation.  See \cref{sec:rankexplained}.
\end{enumerate}

There is also a fourth, more conceptual, payoff that I keep thinking
about: the same operator family $\{\Qabs,\Qref,\Qbist(\kout)\}$
governs what the array \emph{transmits and is absorbed},
\emph{transmits and is scattered to a chosen direction}, and (by the
Kirchhoff dual of \texttt{q\_complement.tex} \S\,6) what the array
\emph{receives as thermal pickup from the same body}.  Three
physical observables, one operator family.  This is the body's full
electromagnetic signature, factored.

\subsection{Voicing the RIS analogy honestly}

I will write \emph{reflective intelligent body} (RIB) for the body
treated as an analytic scatterer in the propagation graph.  The
analogy with RIS is intentional but it is not literal.  A real RIS
has electronically tunable per-element phase shifters; a body does
not.  The right analogy is closer to a \emph{passive,
self-tuning} RIS: the body's surface phase pattern
$e^{-ik_0\kout\cdot\rr}$ is fixed up to pose, and gait
\emph{is} the only knob.  At a typical \SI{100}{ms} pose-update
cadence the twin sees the body as a slowly-tunable RIS with a
$\sim 10$\,Hz refresh rate, which is plenty for most JSAC
applications and is the natural rate the body itself selects by
moving.

This avoids the marketing trap of pretending the body is something
it is not.  The intelligence is in the DTN: the network knows the
body, predicts its scattering, and incorporates it into the precoder
loop.  The body just stands there.

% ============================================================
\section{Where RIB lives in the DTN stack}\label{sec:dtn}
% ============================================================

To make the previous section concrete, here is what RIB actually
adds to the JSAC stack, and what it does not.

\subsection{What is already there}

The JSAC paper's twin has, per served user $u$, the path-space
factorisation $\mathbf{Q}^{(u)}=\Jmat^T\mathbf{M}^{(u)}\Jmat$ with
$\mathbf{M}^{(u)}$ a Hermitian path-kernel evaluated at the body's
location.  $\mathbf{M}^{(u)}$ is built from a per-triangle Fresnel
projection of the path-set $\{\khat_n,\bm\psi_n\}$, integrated over
the body surface with the depth-coupling weight that gives absorbed
power.  This object is the heart of the body twin.

\subsection{What RIB adds, mechanically}

The bistatic body matrix $\Frib^{(n)}(\kout)$ defined in
\cref{def:Fb} below is a \emph{re-weighted surface integral over
the same triangles} with three changes: (i) replace the Fresnel
transmission operator $\Fn(\rr)$ with its reflection sibling
$\Rn(\rr)$ from \texttt{q\_complement.tex}; (ii) drop the
depth-coupling factor $\sqrt{\sigma/(4\alpha_n)}$, since the body
does not absorb the reflected beam; (iii) add a
direction-resolved phase $e^{-ik_0\kout\cdot\rr}$.  None of these
changes the iteration over triangles; they change the per-triangle
weight.  Concretely, the twin's surface loop builds a tensor
\[
  \mathbf{M}^{(u)}_{\mathrm{ext}}(\kout)
  = \int_\Sigma \mathbf{R}(\rr)^H\Pi^{\mathrm{out}}\mathbf{R}(\rr)\,
                e^{-ik_0\kout\cdot\rr}\,\diff A,
\]
which the precoder side queries as a quadratic form
$\mathbf{x}^H\,\Jmat^T\mathbf{M}^{(u)}_{\mathrm{ext}}(\kout)\,\Jmat\,\mathbf{x}$.

\subsection{What it does not add}

RIB is silent on the served-user channel.  $\hh$ is what reciprocity
gives, and reciprocity already includes any body-mediated paths.
The first thing a reviewer will ask is ``what does RIB tell the
served-user channel that pilots do not'', and the honest answer is
``nothing direct''.  RIB earns its keep elsewhere: at the bystander,
in the rank explanation, and in the Pareto bound.

\subsection{The calibration story}\label{sec:calib}

The first-principles RIB pattern assumes a smooth pose-tracked
surface with homogeneous tissue.  Real bodies are clothed, hairy,
and only approximately at the SMPL-X estimate.  RIS deployments deal
with the same gap by treating the analytic model as a structural
prior and absorbing residual amplitude/phase via a per-element
calibration coefficient $\beta$, fit from CSI.  The JSAC paper
(\S\,IV.D) already does exactly this for the path dictionary, with
ridge least-squares against an in-silico amplitude prior.

For RIB the natural extension is a \emph{per-body, per-direction-bin}
calibration coefficient $\beta_b^{(u)}(\Omega)$ that multiplies
$\Frib^{(u)}(\kout)$ inside the bin $\Omega$.  Bin sizes follow the
specular-lobe width $\Delta\theta_{\mathrm{spec}}\sim
\sqrt{\lambda/(2\pi R)}$ from \cref{eq:lobe}, giving roughly
$\sim 100{-}500$ angular bins per body at \SI{28}{GHz}.
Calibration data come from passive observation of the served-user
multipath: subtract the predicted body-mediated paths, and the
residual amplitude per bin gives $\beta_b^{(u)}$.  This is the
RT-calibrated DT pattern: structure from physics, residuals from
data.  Specifically:
\[
   \min_{\bm\beta_b^{(u)}}\;
   \sum_{m\in\mathrm{CSI}}\bigl|h_m
       -\hat h_m^{\mathrm{pred}}(\bm\beta_b^{(u)})\bigr|^2
   +\lambda\,\|\bm\beta_b^{(u)}-\bm 1\|^2,
\]
ridge-regularised around the analytic prior, identical to the JSAC
\S\,IV.D recipe.

The fidelity floor expected from this calibration: $\sim 1$\,dB
on amplitude (matching the JSAC CSI-calibration result for the path
dictionary), $\sim 5{-}10^\circ$ on specular peak position (set by
pose-update granularity), and $\sim 30\%$ on absolute peak amplitude
when the body is fully clothed (set by textile roughness, the
microfacet-BRDF perturbation that the calibration absorbs).  None of
these are paper-killing for the DTN application.

% ============================================================
\section{Setup and assumptions}\label{sec:setup}
% ============================================================

The body is a closed orientable surface
$\Sigma\subset\mathbb{R}^3$ with outward unit normal $\nhat(\rr)$ on
$\rr\in\Sigma$.  The tissue is a homogeneous lossy half-space with
complex refractive index $\tilde n(f)$; layered tissue extensions
follow the same scaffolding with the layered Fresnel coefficients of
Part~I, \S\,3.  The array radiates a precoder
$\mathbf{x}\in\Complex^M$, producing $N$ propagation paths
$\{j(n),\khat_n,\bm\psi_n\}_{n=1}^N$ at the body location, with the
same notation as Part~III.

\begin{enumerate}[label=(A\arabic*),leftmargin=2.4em,itemsep=2pt]
\item\label{a:po} \emph{High-frequency regime.}
Curvature radius of $\Sigma$ exceeds the wavelength,
$\kappa^{-1}\gtrsim\lambda/2\pi$.  At \SI{28}{GHz} this is
$\lambda\approx\SI{1.07}{cm}$; arms, torso, head all comply outside
sharp features (eyebrow ridges, fingers).
\item\label{a:far} \emph{Far-field observer.}  The receiver lives at
$|\mathbf{R}|\gg D^2/\lambda$ from the body, where $D$ is the body
diameter ($D\approx\SI{1.8}{m}$, $D^2/\lambda\approx\SI{300}{m}$ at
\SI{28}{GHz}).  Below this we are in the near-field regime; the
Kirchhoff integral still applies but $1/R$ is replaced by a Fresnel
kernel.  See \cref{sec:nearfield}.
\item\label{a:single} \emph{Single bounce.}  Power that reflects off
one body part and hits another is treated as a higher-order
correction (\texttt{q\_complement.tex} \S\,8.1).  At mmWave
self-visibility is low, so single-bounce captures $\gtrsim 95\%$ of
the scattered power.
\item\label{a:pb} \emph{Pseudo-Brewster regime.}  Per
\texttt{q\_complement.tex} \S\,4 we work where $T_s\approx T_p\approx
T_0$ across $\theta\in[0,\SI{80}{\degree}]$.  This is the mmWave
skin/water regime; sub-\SI{6}{GHz} corrections are
polarisation-dependent and treated separately.
\end{enumerate}

\begin{remark}
The PO/Kirchhoff approximation underpinning everything below is the
same one that converts Mie series into Fresnel-times-projected-area
in Part~I.  The validation route is the same: residual analysis
against full-wave Mie or FDTD.  The TAP submission already shows the
PO approximation buys $\le$\,$5\%$ on absorbed power; we expect the
same on scattered power, with the residual concentrated at
oblique-and-curved patches (eyebrow ridges, fingertips).
\end{remark}

% ============================================================
\section{The body scattering function}\label{sec:Fb}
% ============================================================

\subsection{Reflected surface field}

For a single plane wave $(\khat,\bm\psi)$ incident on the front face,
the Fresnel-reflected field on $\Sigma^+$ is
\begin{equation}\label{eq:Erefl-pw}
  \EE^{\mathrm{re}}(\rr) = \mathbf{R}(\rr;\khat)\,\bm\psi\;
                          e^{-ik_0\khat\cdot\rr}\,H(\mu(\rr;\khat)),
\end{equation}
with $\mu=\nhat\cdot(-\khat)$ and $\mathbf{R}$ the Fresnel reflection
operator from \texttt{q\_complement.tex} Definition~2.1.  Under
\ref{a:pb} this collapses to
\begin{equation}\label{eq:R-collapse}
  \mathbf{R}(\rr;\khat)\approx
   \sqrt{1-T_0}\,e^{i\phi_r(\theta)}\,\Pi(\rr;\khat),
   \qquad \Pi=\ident_3-\khat\khat^{\!T},
\end{equation}
i.e.\ the reflection operator is, locally, a polarisation projector
times a body-universal scalar.  The scalar
$r_0\equiv\sqrt{1-T_0}\,e^{i\phi_r(\theta)}$ depends on local
incidence angle but not on polarisation; the projector $\Pi$ is a
geometric object.

For the array's full superposition the surface-reflected field is
\begin{equation}\label{eq:Erefl-array}
  \EE^{\mathrm{re}}(\rr)=
  \sum_n \mathbf{R}(\rr;\khat_n)\,\bm\psi_n\,
        e^{-ik_0\khat_n\cdot\rr}\,H(\mu_n)\;\cdot\,x_{j(n)},
\end{equation}
which equals $\mathbf{G}^{\mathrm{re}}(\rr)\mathbf{x}$ with
$\mathbf{G}^{\mathrm{re}}$ the reflection-channel matrix introduced in
\texttt{q\_complement.tex} Eq.~(7).

\subsection{Kirchhoff far-field integral}

Standard physical optics gives the far-zone scattered field at
$\mathbf{R}=R\,\rhat$ in direction $\rhat=\kout$ as the surface
integral of the equivalent currents.  Under \ref{a:po}, the magnetic
current dominates at smooth specular patches; the scattered electric
field reads
\begin{equation}\label{eq:KH}
  \EE^{\mathrm{sc}}(\mathbf{R})
  =\frac{ik_0\,e^{ik_0 R}}{4\pi R}\,
   \Pi^{\mathrm{out}}\!\!\int_\Sigma
   \EE^{\mathrm{re}}(\rr)\,e^{-ik_0\kout\cdot\rr}\,
   H(\nhat\cdot\kout)\,\diff A(\rr) + O(R^{-2}),
\end{equation}
where $\Pi^{\mathrm{out}}=\ident_3-\kout\kout^{\!T}$ enforces the
transversality of radiation in the outgoing direction, and the second
Heaviside selects the lit patches visible to the observer.  The
factor $ik_0/(4\pi R)$ is the standard Stratton--Chu far-field
normalisation in Gaussian-vector form, with the $1/(4\pi)$ inherited
from the 3D scalar Green's function.

Substituting \cref{eq:Erefl-array}:

\begin{definition}[Body scattering matrix]\label{def:Fb}
For incident path $n$ with direction $\khat_n$ and polarisation
$\bm\psi_n$, define the per-path body scattering matrix at outgoing
direction $\kout$ by
\begin{empheq}[box=\tcbhighmath]{equation}\label{eq:Fb}
  \Frib^{(n)}(\kout)=
  \frac{ik_0}{4\pi}\,
  \Pi^{\mathrm{out}}\!\!\int_\Sigma\!\!
  \mathbf{R}(\rr;\khat_n)\,
  e^{-ik_0(\kout-\khat_n)\cdot\rr}\,
  H(\mu_n(\rr))\,H(\nhat\cdot\kout)\,\diff A.
\end{empheq}
The full body scattering matrix at outgoing direction $\kout$ for
array precoder $\mathbf{x}$ is
\begin{equation}\label{eq:Fb-x}
  \EE^{\mathrm{sc}}(R\kout)\,R\,e^{-ik_0R}=
  \sum_n \Frib^{(n)}(\kout)\,\bm\psi_n\,x_{j(n)}
  =\Frib(\kout)\,\mathbf{x}\in\Complex^3,
\end{equation}
with $\Frib(\kout)=\sum_n\Frib^{(n)}(\kout)\bm\psi_n\mathbf{e}_{j(n)}^{\!T}
\in\Complex^{3\times M}$.
\end{definition}

The integrand has a transparent structure: a stationary-phase factor
$e^{-ik_0(\kout-\khat_n)\cdot\rr}$, a polarisation-and-amplitude
filter $\mathbf{R}(\rr;\khat_n)$, and two visibility gates.  The
stationary-phase points are exactly the specular patches:
$\rr^*$ where the local normal bisects $\khat_n$ and $\kout$,
$\nhat(\rr^*)\propto \kout-\khat_n$.

\subsection{Far-field intensity and bistatic RCS}

The far-field power flux is
$\Sref^{\mathrm{ff}}(R\kout)=\|\Frib(\kout)\mathbf{x}\|^2/(Z_0 R^2)$.
Comparing with the conventional bistatic-RCS definition
\begin{equation}\label{eq:rcs-def}
  \sigb(\kin,\kout;\bm\psi)=\lim_{R\to\infty}
  4\pi R^2\,\frac{|\EE^{\mathrm{sc}}(R\kout)|^2}{|\EE^{\mathrm{inc}}|^2},
\end{equation}
yields, for a single incident plane wave $(\kin=\khat,\bm\psi)$ of
unit amplitude,
\begin{equation}\label{eq:RCS-PO}
  \boxed{\sigb(\kin,\kout;\bm\psi)=
  \frac{k_0^2}{\pi}\,
  \bigl|\Pi^{\mathrm{out}}\!\!\int_\Sigma\!\!
  \mathbf{R}(\rr;\kin)\bm\psi\,
  e^{-ik_0(\kout-\kin)\cdot\rr}
  H(\mu)H(\nhat\cdot\kout)\,\diff A\bigr|^2.}
\end{equation}
This is the body's bistatic RCS in physical optics, parametrically
identical to the rough-rigid-body PO formulae of textbook radar but
with the Fresnel reflection operator in place of a perfect-conductor
unity factor.

\begin{lemma}[Specular peak height]\label{lem:peak}
Under \ref{a:pb} and the stationary-phase approximation around a
single specular patch of curvature radii $(R_1,R_2)$,
\begin{equation}\label{eq:peak}
  \sigb(\kin,\kout^*;\bm\psi)\big|_{\mathrm{spec}}
  \approx \pi\,(1-T_0)\,R_1 R_2\,|\bm\psi|^2,
\end{equation}
where $\kout^*$ is the specular direction and the body universal
constant is the same $1-T_0$ as in \texttt{q\_complement.tex}.
\end{lemma}

Sketch.  The standard PO result $\sigma=\pi R_1 R_2$ for a perfectly
conducting smooth scatterer at specular is multiplied by the local
power reflectance $|r_0|^2=1-T_0$ when the boundary condition is
soft (Fresnel) instead of perfect.  Polarisation averaging is taken
care of by the projector identity \cref{eq:R-collapse}.

\Cref{lem:peak} gives the right ballpark: a torso of $R\approx
\SI{0.15}{m}$ produces a specular peak of order
$\pi\cdot 0.37\cdot R^2\approx \SI{2.6e-2}{m^2}$ or about
$-15.8\,\mathrm{dBsm}$ at \SI{28}{GHz}, consistent with the
$0\,$dBsm placeholder used in the JSAC paper for body-scale
mono-static radar (which absorbs torso-plus-head plus a bit of
multipath).

% ============================================================
\section{The directional exposure operator}\label{sec:Qbist}
% ============================================================

The single-direction analog of the operator-level book-keeping in
\texttt{q\_complement.tex} is now obvious.

\begin{definition}[Directional exposure operator]\label{def:Qbist}
For each outgoing direction $\kout\in S^2$, define the Hermitian
$M\!\times\!M$ matrix
\begin{equation}\label{eq:Qbist}
  \Qbist(\kout)\equiv
  \frac{1}{Z_0}\,\Frib(\kout)^H\,\Frib(\kout).
\end{equation}
Then $\mathbf{x}^H\Qbist(\kout)\mathbf{x}$ is the body's
far-field power flux per unit solid angle into direction $\kout$,
multiplied by $R^2/(\text{nothing extra})$.
\end{definition}

This is the angular density of the body's outgoing radiation, but it
is \emph{not} simply the angular density of $\Qref$ until we
separate the back- from the forward-scatter.  This is a real subtlety
worth getting right.

\subsection{Back-scatter, forward-scatter, and what each operator counts}

The Kirchhoff far-field integral \cref{eq:Fb} contains
\emph{both} the specular/back-scatter contribution (which is what
$\Qref$ from \texttt{q\_complement.tex} Definition~2.2 counts) and a
forward-scatter contribution (the Babinet shadow lobe at
$\kout\approx\kin$, of integrated cross-section
$\sim A_{\mathrm{proj}}$).  The forward lobe is the diffraction-induced
forward field that interferes destructively with $\EE^{\mathrm{inc}}$
to produce the geometric shadow behind the body.  $\Qref$ in
\texttt{q\_complement.tex} only books power that
\emph{leaves the body in a non-forward direction}; the forward
shadow is paired with the absorbed and missed channels in the
Stratton--Chu energy budget, not the reflected one.

Concretely the angular bookkeeping splits as
\begin{equation}\label{eq:split}
  \int_{S^2}\!\!\Qbist(\kout)\diff\Omega
  =\underbrace{\int_{\Omega_{\mathrm{back}}}\!\!\Qbist(\kout)\diff\Omega}_{\approx\,\Qref}
  +\underbrace{\int_{\Omega_{\mathrm{fwd}}}\!\!\Qbist(\kout)\diff\Omega}_{\approx\,A_{\mathrm{proj}}\mathbf{I}/(MZ_0)},
\end{equation}
where $\Omega_{\mathrm{fwd}}$ is the small forward cone of solid
angle $\sim\lambda^2/A_{\mathrm{proj}}$ around $\kout=\kin$ and
$\Omega_{\mathrm{back}}$ is everything else.

\begin{theorem}[Bookkeeping reconciliation, corrected]\label{thm:reconcile}
Under \ref{a:po}--\ref{a:pb}, the back-hemisphere integral of the
directional operator equals the reflection operator of
\texttt{q\_complement.tex}:
\begin{empheq}[box=\tcbhighmath]{equation}\label{eq:reconcile}
  \int_{\Omega_{\mathrm{back}}}\!\!\Qbist(\kout)\,\diff\Omega(\kout)
  \;=\;\Qref + O(\bar\alpha^2),
\end{empheq}
with the residual being the same multi-bounce correction as in
\texttt{q\_complement.tex} Eq.~(20).  The forward-cone contribution
\emph{adds} to the missed and absorbed channels rather than the
reflected one, restoring the global identity
$\Qabs+\Qref+\Qmis=\mathbf{I}$ of \texttt{q\_complement.tex}
Theorem~3.1.
\end{theorem}

Sketch.  Substitute \cref{eq:Fb}, expand the $|\cdot|^2$ in
\cref{eq:RCS-PO}, integrate over $\kout$, and use Parseval on the
phase factor $e^{-ik_0\kout\cdot\rr}$ to collapse the
double-surface integral to the single-surface integral that defines
$\Qref$.  The Parseval step is exact only on the angular subset that
excludes the forward stationary point $\kout=\kin$, where the phase
$e^{-ik_0(\kout-\kin)\cdot\rr}$ degenerates and the shadow lobe
appears.  Restricting integration to $\Omega_{\mathrm{back}}$ removes
that subset; the residual is the same self-visibility cascade as
\texttt{q\_complement.tex} Remark~2.5.  The forward contribution is
its own object and needs the Stratton--Chu paired-current treatment
to be partitioned into $\Qabs$ and $\Qmis$; see
\cref{rem:forwardshadow}.

\begin{remark}[The forward-shadow lobe in practice]\label{rem:forwardshadow}
For the demo of \cref{sec:numerical}, naively integrating
$|\Frib|^2$ over the full sphere gives a total scattered power
$\sim (2-T_0)\cdot A_{\mathrm{proj}}/Z_0$, the
Babinet-doubled cross-section, not $(1-T_0)\cdot A_{\mathrm{proj}}/Z_0$.
Restricting to $\Omega_{\mathrm{back}}$, or equivalently subtracting
the forward lobe of solid angle $\sim\lambda^2/A_{\mathrm{proj}}$,
recovers the operator-level number to within the PO residual.  The
companion script \texttt{rib\_far\_field\_demo.py} exposes both
totals and reports the shadow-corrected value as the apples-to-apples
comparison with $\mathbf{x}^H\Qref\mathbf{x}$.
\end{remark}

\begin{corollary}[Trace identity]\label{cor:trace}
The body's monostatic radar cross-section, integrated over all
incident directions and weighted by the array, equals the trace of
$\Qref$:
\[
  \int_{S^2}\!\!\sigb(\kin,\kin)\,\diff\Omega(\kin)
  \cdot \tfrac{1}{4\pi}=\frac{\tr\Qref}{M}\cdot 4\pi,
\]
modulo the same multi-bounce residual.  In words: the body's
isotropic-monostatic average RCS is, up to a body-universal prefactor,
the body's mean albedo from \texttt{q\_complement.tex} \S\,7.
\end{corollary}

\begin{remark}[Why ``intelligent'']
$\Qbist(\kout)$ is to RIB what the cascaded-channel matrix
$\mathbf{H}\mathbf{\Phi}\mathbf{G}$ is to RIS, with $\mathbf{\Phi}$
replaced by a per-direction \emph{passive} Hermitian operator the
array can compute from pose and tissue alone.  The body cannot tune
its scattering, but the array can predict it.  In the propagation
graph the body acts as an analytic node.
\end{remark}

% ============================================================
\section{Specular versus diffuse}\label{sec:specdiff}
% ============================================================

A body has smooth-on-the-wavelength patches (torso) and
roughness-on-the-wavelength features (hair, fabric, fingers).
The Kirchhoff integral handles both, but its structure looks different
in the two regimes.

\subsection{The smooth patch}

On a smooth patch of dimensions $\gg\lambda$ the stationary-phase
approximation isolates a single specular point with
$\nhat\propto\kout-\kin$.  The Hessian of the phase at that point is
proportional to the local Gaussian curvature, giving
\cref{lem:peak}.  The full $\sigb(\kin,\kout)$ is then a narrow lobe
of angular width
\begin{equation}\label{eq:lobe}
  \Delta\theta_{\mathrm{spec}}\sim
  \sqrt{\frac{\lambda}{2\pi R_{\min}}},
\end{equation}
with $R_{\min}=\min(R_1,R_2)$ the smaller principal curvature radius.
For a torso ($R_{\min}\approx\SI{15}{cm}$) at \SI{28}{GHz} this is
$\approx \SI{2.4}{\degree}$.  For a head
($R_{\min}\approx\SI{8}{cm}$) it is $\approx \SI{3.3}{\degree}$.  The
specular reflection is forward-peaked enough that ``a person standing
in a beam'' produces a tight, predictable secondary lobe in the
specular direction at mmWave.

\subsection{The diffuse part}

Outside the specular cone, the body scatters diffusely.  In PO this
diffuse contribution comes from sub-wavelength surface features (the
mesh's high-frequency content, plus textile roughness which is not
in the mesh).  A standard separation writes
\begin{equation}\label{eq:rho-diffuse}
  \sigb(\kin,\kout)
  =\sigb^{\mathrm{spec}}(\kin,\kout)+\sigb^{\mathrm{diff}}(\kin,\kout),
\end{equation}
with the diffuse part well approximated by Lambert's law times the
albedo-density factor $1-T_0$:
\begin{equation}\label{eq:Lambert}
  \sigb^{\mathrm{diff}}(\kin,\kout)\approx
  \frac{(1-T_0)\,\mu^{\mathrm{in}}\,\mu^{\mathrm{out}}}{\pi}
  \,A_{\mathrm{lit}}(\kin,\kout),
\end{equation}
where $A_{\mathrm{lit}}$ is the area lit by both
$\kin$ and $\kout$.  The constant $1-T_0\approx 0.37$ is the body's
albedo; the cosine factors are the Lambertian foreshortening; the
joint visibility comes from \cref{eq:Fb}.  This is the same model
used in computer graphics for human skin BRDFs and gives a Sionna-style
``diffuse material'' definition for a body.

For a torso at \SI{28}{GHz} with body diameter
$D\approx\SI{0.4}{m}$, the diffuse component carries
$\sim 60\%$ of the total scattered power, and the specular lobe
carries $\sim 40\%$, with $\Delta\theta_{\mathrm{spec}}$ as above.
This is a conjecture I want a numerical check on; see
\cref{sec:numerical}.

% ============================================================
\section{Per-direction pseudo-Brewster locking}\label{sec:lock}
% ============================================================

The eigenvector locking of \texttt{q\_complement.tex} Prop~4.2 was a
statement about the integrated Hermitian operators
$\Qref\approx\frac{1-T_0}{T_0}\Qabs$.  The directional version is
slightly more delicate, and slightly more interesting.

\subsection{The directional absorbed-exposure operator}

For an outgoing direction $\kout$, define the
\emph{$\kout$-conjugate} exposure operator
\begin{equation}\label{eq:Qabs-kout}
  \Qabs^{(\kout)}=\!\!\int_\Sigma\!\!
  \Gtil(\rr)^H\Gtil(\rr)\,
  H(\nhat(\rr)\cdot\kout)\,\diff A(\rr).
\end{equation}
This is exactly the standard exposure operator restricted to the
patches that are visible to a far-field observer in direction $\kout$.

\begin{proposition}[Per-direction locking]\label{prop:perdir}
Under \ref{a:po}--\ref{a:pb},
\begin{equation}\label{eq:perdir}
  \Qbist(\kout)\;\approx\;
  \frac{1-T_0}{T_0}\,\Qabs^{(\kout)} + \Delta(\kout),
\end{equation}
with $\|\Delta(\kout)\|/\|\Qabs^{(\kout)}\|=O(5\%)$ uniformly in
$\kout$.
\end{proposition}

Sketch.  Repeat the proof of \texttt{q\_complement.tex} Prop~4.2 with
the extra outgoing-visibility gate $H(\nhat\cdot\kout)$ on both
$\Qbist$ and $\Qabs^{(\kout)}$.  The pseudo-Brewster collapse
\cref{eq:R-collapse} carries over verbatim because it is local: at
each visible point $\rr$, the polarisation projector $\Pi$ is the
same on both sides, the only difference is the scalar
$T_0\leftrightarrow 1-T_0$ and the outgoing phase
$e^{ik_0\kout\cdot\rr}$.  The phase is unitary in the surface
integral that builds the Gram, so cancels.  $\Pi^{\mathrm{out}}$ on
$\Qbist$ is a unitary-on-the-output projector and it commutes with
the surface integral (its action is on the polarisation index, not
the spatial index).

\begin{corollary}[Sensing-exposure duality, angular form]
For every outgoing direction $\kout$, the precoder that maximises
backscatter into $\kout$ is, to within $5\%$, the precoder that
maximises absorbed power on the $\kout$-visible patches of the body.
\end{corollary}

This is the angular refinement of the classic statement.  The integral
version was: ``the worst beam for the bystander is the best beam for a
monostatic radar.''  The angular version is: ``\emph{for any
receiver placement}, the worst beam is the best beam.''  Place the
receiver behind the body and the bound applies to the scattered cone
out the back.  Place it overhead and the bound applies to the cone
out the top.  This is what makes RIB a bona fide bistatic-sensing
construct rather than a monostatic one.

% ============================================================
\section{The capacity-exposure Pareto bound}\label{sec:pareto}
% ============================================================

The single most useful consequence of per-direction locking is that
the body-network's contribution to channel rank cannot be made large
without making exposure on the bodies large.  In a DTN that already
budgets exposure, this is the kind of bound that turns into an
allocation rule.

\subsection{The capacity gain object}

Let the array serve $K$ users with channels
$\hh_1,\ldots,\hh_K$ and use a precoder
$\mathbf{X}=[\mathbf{x}_1,\ldots,\mathbf{x}_K]$ with sum power
$\|\mathbf{X}\|_F^2\le P_{\mathrm{tx}}$.  Without bodies in the
scene, the signal channel for user $u$ is the $K$ direct paths
$\hh_u^{\mathrm{los}}$.  With $K_b$ bodies in the scene at positions
$\rr_b$, each body contributes additional second-bounce paths to
$\hh_u$ via $\Frib^{(b)}(\hat\kout_{b\to u})$ as in
\cref{eq:bvec-rib}.  The body-network capacity gain is
\begin{equation}\label{eq:Gbody}
  G_{\mathrm{body}}(\mathbf{X})\equiv
  \log\det\!\Bigl(\mathbf{I}+\textstyle\frac{1}{N_0}\sum_u
            \mathbf{X}^H\hh_u^{\mathrm{full}}\hh_u^{\mathrm{full}\,H}\mathbf{X}\Bigr)
  -\log\det\!\Bigl(\mathbf{I}+\textstyle\frac{1}{N_0}\sum_u
            \mathbf{X}^H\hh_u^{\mathrm{los}}\hh_u^{\mathrm{los}\,H}\mathbf{X}\Bigr),
\end{equation}
the difference between the full-channel and LOS-only mutual
information.  It is non-negative and bounded by the rank lift the
body-network provides.

\subsection{The exposure constraint}

The DTN enforces, for each user-or-bystander body $u$,
$\mathbf{x}_u^H\Qabs^{(u)}\mathbf{x}_u\le L_{\mathrm{RL}}^{(u)}$, the
basic-restriction limit.  Define the array-level exposure dose
\begin{equation}\label{eq:dose}
  D(\mathbf{X})\equiv \sum_u
   \mathrm{tr}\bigl(\mathbf{x}_u^H\Qabs^{(u)}\mathbf{x}_u\bigr)
  =\sum_u\mathbf{x}_u^H\Qabs^{(u)}\mathbf{x}_u.
\end{equation}

\subsection{The bound}

\begin{theorem}[Capacity-exposure Pareto bound]\label{thm:pareto}
Under \ref{a:po}--\ref{a:pb}, every body's contribution to
$G_{\mathrm{body}}$ is upper-bounded by a function of that body's
absorbed dose:
\begin{empheq}[box=\tcbhighmath]{equation}\label{eq:pareto}
  G_{\mathrm{body}}(\mathbf{X})\;\le\;
  K_b\,\log\!\Bigl(1+\tfrac{1}{N_0}\,\tfrac{1-T_0}{T_0}\,
                   \mathcal{K}_{\mathrm{rx}}\,
                   D(\mathbf{X})\Bigr)
  +\,O(5\%),
\end{empheq}
where $\mathcal{K}_{\mathrm{rx}}$ collects the receive-aperture
collection factors of the user-side antennas and is independent of
the array's precoder.
\end{theorem}

Sketch.  Bound $G_{\mathrm{body}}$ above by Hadamard's inequality on
$\det$ to get a per-body contribution
$\log(1+\tfrac{1}{N_0}\mathbf{x}_u^H\Qbist^{(u)}(\hat\kout_{b\to u})\mathbf{x}_u\cdot\mathcal{K}_{\mathrm{rx}})$.
Apply \cref{prop:perdir} to replace
$\mathbf{x}_u^H\Qbist^{(u)}(\kout)\mathbf{x}_u$ with
$\frac{1-T_0}{T_0}\mathbf{x}_u^H\Qabs^{(u),(\kout)}\mathbf{x}_u$.
Bound $\Qabs^{(u),(\kout)}\preceq\Qabs^{(u)}$ since the
$\kout$-visibility gate restricts the integration domain.  Sum
over bodies.  $\square$

The structural content of \cref{thm:pareto}: \emph{each watt of
body-mediated capacity costs at least $T_0/(1-T_0)\approx 0.32$\,W
of absorbed dose.}  At mmWave, the Pareto slope is
body-universal because $T_0$ is body-universal under
pseudo-Brewster.

\subsection{What the bound implies for DTN allocation}

The bound is exactly tight when the precoder aligns with the top
eigenvector of $\Qabs^{(u),(\hat\kout_{b\to u})}$, which by
\cref{prop:perdir} is also the top eigenvector of
$\Qbist^{(u)}(\hat\kout_{b\to u})$.  A DTN that wants to maximise
RIB-mediated capacity \emph{should} steer along that eigenvector,
which is the same eigenvector exposure compliance forbids the
precoder from steering along.  Two consequences worth pulling out:

\begin{enumerate}[leftmargin=1.6em,itemsep=2pt]
\item \textbf{Capacity and exposure are not independent axes.}  The
JSAC paper's Pareto frontier between sum-rate and dose budget,
which the multi-user QCQP solver maps out, is the \emph{same}
frontier as the body-network capacity gain plus the LOS sum-rate.
There is one Pareto, not two.
\item \textbf{Body-mediated capacity has a hard ceiling.}  The
maximum rank lift the body-network can ever deliver, summed across
all $K_b$ bodies, scales as $\log\det$ of the diagonal-stacked
$\Qbist^{(u)}(\hat\kout_{b\to u})$.  Under \cref{prop:perdir},
this is bounded by $\frac{1-T_0}{T_0}\sum_u\lambda_1(\Qabs^{(u)})$,
which is the same eigenvalue the basic restriction caps at
$L_{\mathrm{RL}}^{(u)}/P_{\mathrm{tx}}$.  Hard ceiling, body by
body.
\end{enumerate}

The DT-aware allocation rule that drops out: do not waste exposure
budget on body-mediated paths that you can serve via other users'
LOS, because every dose watt buys at most a fixed amount of
body-mediated capacity, set by $(1-T_0)/T_0$.

\begin{remark}[How tight is ``$\le$''?]\label{rem:pareto-tight}
The inequality is loose by Hadamard plus the visibility-restriction
relaxation $\Qabs^{(u),(\kout)}\preceq\Qabs^{(u)}$.  In the regime
where each body contributes a single dominant outgoing direction (a
served user behind the body, or a directional bystander
detection geometry), Hadamard is approximately tight and the
visibility gate is loose by the fraction $A_{\mathrm{vis}}/A_{\mathrm{lit}}$.
For a torso seen from a single direction this ratio is $\sim 0.5$,
so the bound underestimates feasible $G_{\mathrm{body}}$ by about
$3$\,dB.  In dense crowds with many bodies and many outgoing
directions, the bound tightens because Hadamard becomes a closer
characterisation of the per-direction contributions and the
visibility gate fractions average toward
$\Omega_{\mathrm{lit}}/4\pi\approx 1/4$.
\end{remark}

% ============================================================
\section{Where the rank-3-to-4 observation comes from}\label{sec:rankexplained}
% ============================================================

The JSAC paper observes, on plaza-scale Brussels Grand Place
runs, that the per-user $\mathbf{Q}^{(u)}$ has effective rank
$3{-}4$ at $99\%$ trace, and attributes this to
``low-rank generically.''  RIB gives a derivation.

\subsection{Body-side angular-mode count}

The integrated-over-$\kout$ operator $\Qref^{(u)}=
\int\Qbist^{(u)}(\kout)\diff\Omega$ is a Hermitian PSD matrix on
$\Complex^M$.  Its rank is bounded above by the number of independent
\emph{body-side} angular modes the surface supports, which is the
same as the number of independent path arrival directions
$\{\khat_n\}_n$ that hit the lit area (per
\texttt{q\_complement.tex} Prop~5.2),
\begin{equation}\label{eq:Meff}
  M_{\mathrm{eff}}^{(u)}\equiv
  \mathrm{rank}\,\Qref^{(u)}
  \le \min(M,\,N_{\mathrm{paths-on-body}}^{(u)}).
\end{equation}
At \SI{28}{GHz} for a thelonious phantom and a typical
plaza-rooftop URA, $N_{\mathrm{paths-on-body}}\sim 20{-}40$ paths
contribute meaningfully (the rest are negligibly weak), so
$M_{\mathrm{eff}}^{(u)}\sim 10{-}30$ in this geometry.

\subsection{Array-side reachability cap}

The precoder lives in $\Complex^M$.  The array can excite, at a
fixed body location, at most $\dim\mathrm{span}\,\{\Jmat\,\bm a(\khat_n)\}_n$
independent surface fields, where $\bm a$ is the per-element
steering response and the union is over body-incident paths $n$.
Call this the \emph{array reachable rank} $R_{\mathrm{arr}}^{(u)}$.

\subsection{The product}

The exposure operator
$\mathbf{Q}^{(u)}=\Jmat^T\mathbf{M}^{(u)}\Jmat$ is a sandwich, and
its effective rank is the minimum of the two:
\begin{equation}\label{eq:rank-product}
  \mathrm{rank}_{\mathrm{eff}}\,\mathbf{Q}^{(u)}
   \;\approx\;\min\bigl(M_{\mathrm{eff}}^{(u)},\,R_{\mathrm{arr}}^{(u)}\bigr).
\end{equation}
For the JSAC paper's plaza scenario ($M=64$ panel, single body at
$\sim\SI{50}{m}$, $\sim 5^\circ$ angular spread of body-incident
paths), $R_{\mathrm{arr}}^{(u)}\sim 3{-}5$.  The product is $\sim 4$,
matching the empirical effective rank.

\subsection{Why this is a useful explanation}

\Cref{eq:rank-product} predicts how the rank should scale with
parameters that the JSAC paper varies:
\begin{itemize}[leftmargin=1.6em,itemsep=2pt]
\item More array elements $M$, fixed body distance: $R_{\mathrm{arr}}$
grows roughly as $\sqrt{M}$ (angular resolution improves as
$1/\sqrt{M}$); rank should increase.  This is what the JSAC paper's
$M=256$ ablation shows.
\item Closer body, fixed array: $M_{\mathrm{eff}}$ grows because
the body subtends more solid angle from the array; rank should
increase.  Predicts that near-field bodies have higher
$\mathrm{rank}\,\mathbf{Q}$.
\item Multiple bodies: each contributes its own
$M_{\mathrm{eff}}^{(u)}$ subspace.  These subspaces are partially
overlapping (paths to nearby bodies share array-side directions)
but the union has dimension $\sim K_b\,M_{\mathrm{eff}}$ in the
sparse-crowd limit.  This is the rank source for
\cref{sec:multi}'s capacity claim and matches the JSAC paper's
observation that the multi-user trace grows sub-linearly with $K$.
\end{itemize}

In short: the empirical rank observation is a derivable consequence
of the bistatic-mode count of the body, intersected with the
angular-resolution capacity of the array.  Neither factor required
any fitting or measurement.

% ============================================================
\section{Second-bounce ray tracing}\label{sec:secondbounce}
% ============================================================

Concretely: how does an existing Sionna ray tracer ingest RIB?

\subsection{What Sionna sees}

Per \texttt{src/aegis/modal\_rt/sionna\_tracer.py}, paths come back
from Sionna with $(j(n),\khat_n,\bm\psi_n)$ at any spatial probe.
A second-bounce path TX$\to$body$_u\to$RX is constructed in three
stages.

\begin{enumerate}[leftmargin=1.6em,itemsep=2pt]
\item \emph{Forward leg.}  Probe the array's path set at the body's
centre $\rr_u$.  This gives the per-element incidence
$\{\khat_n,\bm\psi_n\}$ and propagation phase to body $u$.  Already
in the dictionary; cost zero.
\item \emph{Body scattering.}  Apply the per-direction operator
$\Qbist(\kout)$ for the outgoing direction
$\kout=(\rr_{\mathrm{RX}}-\rr_u)/|\cdot|$.  This is a single Hermitian
quadratic form against $\mathbf{x}$ at the precoder evaluation step.
The operator itself is precomputed once per body+frequency from the
mesh and tissue, by a routine that costs the same as one $\Qabs$
build per outgoing direction or one Kirchhoff integral evaluation.
\item \emph{Backward leg.}  Free-space propagation from $\rr_u$ to
$\rr_{\mathrm{RX}}$, including any further reflections in Sionna's
scene (walls, ground).  Sionna handles this once we feed the body's
scattered field as a new emitting point source with anisotropic
pattern $\Frib(\kout)\mathbf{x}$ and origin $\rr_u$.
\end{enumerate}

The body becomes a \emph{virtual transmitter} at $\rr_u$ whose
emission pattern is computed from its own pose and the array's
precoder.  In Sionna this is a custom material plus an envelope
``re-emit with attached pattern'' wrapper; both are existing primitive
operations.

\subsection{Cost}

For one body and a fixed precoder, computing $\Frib(\kout)$ across a
sphere of $L\times L$ outgoing directions costs
$O(M_\triangle\cdot L^2)$ where $M_\triangle$ is the number of
front-facing triangles.  For \SI{28}{GHz} and Thelonious,
$M_\triangle\sim 10^4$ and $L\sim 64$ gives $\sim 4\cdot 10^7$
complex multiplications, well under a second on CPU and microseconds
on JAX-vmap.  For $K$ bodies in a scene this is $K\times$ that.  The
cost is dominated by the operator build, not the per-precoder
evaluation; the latter is one quadratic form.

\subsection{What changes for the receiver}

If the receiver is a UE that also has its own body channel, the second
bounce now contributes to $\hh$.  In the JSAC notation
$\hh=\Jmat^T\bvec$, the $\bvec$ vector gets an extra term
$\bvec_{u\to\mathrm{UE}}$ collecting all paths that go through body
$u$ before reaching the UE.  Concretely:
\begin{equation}\label{eq:bvec-rib}
  \bvec\to \bvec + \sum_{u} \bvec_{u\to\mathrm{UE}},
  \qquad
  \bvec_{u\to\mathrm{UE}}=
  \frac{e^{ik_0|\rr_{\mathrm{UE}}-\rr_u|}}{4\pi|\rr_{\mathrm{UE}}-\rr_u|}\,
  \mathbf{C}_R(\hat{\bm{k}}_{u\to\mathrm{UE}})^{\!H}\,\Frib^{(u)}(\hat{\bm{k}}_{u\to\mathrm{UE}})^{\!H}\,
  \bm\psi_{u}^{\!*}\cdots,
\end{equation}
where $\mathbf{C}_R$ is the receive antenna projection of the
JSAC paper.  This is just a longer phase chain; the per-element
structure is identical to a regular wall reflection in Sionna, modulo
the body-scattering operator in the middle.

% ============================================================
\section{Multi-body RIB networks}\label{sec:multi}
% ============================================================

In a plaza with $K$ bodies the propagation graph is enriched by
$K\cdot(K+1)$ additional links: each body becomes a virtual emitter
toward each UE and toward each other body.  The distinct objects to
track are:

\begin{itemize}[leftmargin=1.6em,itemsep=2pt]
\item per-body bistatic operators $\Qbist^{(u)}(\kout)$;
\item body-to-body cross-operators $\Qbist^{(u,v)}$ that capture the
power leaving body $u$ and arriving at body $v$;
\item the global Gram of the cascade
$\sum_u\Frib^{(u)}\to\sum_{u,v}\Frib^{(v)}\Frib^{(u)}+\cdots$.
\end{itemize}

The cascade is a Neumann series in the body-to-body visibility kernel
$\mathcal{K}_{\mathrm{bb}}$ with $\|\mathcal{K}_{\mathrm{bb}}\|\le
(1-T_0)\bar V_{\mathrm{cross}}$, where $\bar V_{\mathrm{cross}}$ is
the average cross-body visibility.  In a plaza of $K$ standing
bystanders separated by $\sim\SI{1}{m}$ at $\sim\SI{1.7}{m}$ tall,
$\bar V_{\mathrm{cross}}\sim 0.05$, so the series converges with
contraction $\sim 0.02$ per term and second-and-higher bounces
contribute $\le 5\%$.  The leading term is the single-bounce
$\sum_u\Qbist^{(u)}$; the cross-term is $K(K-1)$ small contributions.

\subsection{Distributed reflector array}

The vector picture is striking.  Each body adds an effective row to a
``passive RIS'' of size $K\cdot M_{\mathrm{eff}}$, where
$M_{\mathrm{eff}}$ is the number of independent angular modes a body
supports (analogous to the rank of $\Qbist^{(u)}(\kout)$ averaged
over $\kout$).  At \SI{28}{GHz} a body has roughly
$M_{\mathrm{eff}}\sim 10$ to $30$ independent angular modes (this is
a conjecture; numerical measurement of $\mathrm{rank}\Qbist(\kout)$
across $\kout$ for a few configurations is a one-day experiment).  In
a $K=50$ plaza this lifts the channel rank by a factor of order
$50\cdot 20=10^3$ over the no-bodies baseline.  In dense crowds the
RIB enhancement of MIMO capacity could be substantial.  This is what
the JSAC paper observes phenomenologically as ``effective rank
$3{-}4$ at $99\%$ trace'' but does not attribute to body multipath.

\subsection{The exposure-for-free claim}

The key bookkeeping observation is that the RIB pipeline computes
$\Qbist^{(u)}(\kout)$ at the same surface-field evaluation points
that define $\Qabs^{(u)}$.  Concretely the ray-tracer-side cost is

\begin{tabular}{ll}
\textit{Sionna w/o bodies:} & TX$\to$walls$\to$RX paths only.\\
\textit{Sionna with bodies as blockers:} & paths through bodies dropped.\\
\textit{Sionna with RIB:} & paths through bodies emit secondary lobes,\\
& computed by an extra body-Kirchhoff integral.\\
& The integral is dominated by surface-field evaluation\\
& at body triangles, exactly what $\Qabs$ already needs.
\end{tabular}

So a JSAC-style scene that wants both sensing-via-bodies and
exposure-aware MIMO pays the same compute budget as exposure-aware
MIMO alone.  The sensing channel is a free by-product.

% ============================================================
\section{Near-field corrections}\label{sec:nearfield}
% ============================================================

When the receiver lives within the body's Fraunhofer distance
$D^2/\lambda$, the $1/R$ in \cref{eq:KH} is replaced by the
Fresnel-Kirchhoff kernel
$e^{ik_0|\rr-\mathbf{R}|}/|\rr-\mathbf{R}|$.  The body scattering
matrix becomes
\begin{equation}\label{eq:Fb-NF}
  \Frib^{(n)}(\mathbf{R})=
  \frac{ik_0}{4\pi}\,\Pi^{\mathrm{out}}(\mathbf{R})\!
  \int_\Sigma\!\!\mathbf{R}(\rr;\khat_n)
  \frac{e^{ik_0(|\rr-\mathbf{R}|-\khat_n\cdot\rr)}}{|\rr-\mathbf{R}|}
  \cdot\,(\text{visibility})\,\diff A.
\end{equation}
For receivers $\gtrsim D$ away from the body, second-order phase
expansion of $|\rr-\mathbf{R}|$ around the body centroid recovers a
Fresnel-corrected far-field with an additional quadratic phase term;
this is the same structure as the near-field exposure operator of
Paper~C \S\,3.2.2.  The path-space factorisation $\mathbf{Q}=
\mathbf{J}^T\mathbf{M}\mathbf{J}$ persists; the kernel $\mathbf{M}$
acquires the Fresnel quadratic phase.  Practically, near-field RIB is
relevant for handheld devices (UE in the user's hand) and for
in-bus/in-train scenarios; both are deployment cases the JSAC paper
acknowledges as future work.

% ============================================================
\section{Special cases worth checking}\label{sec:special}
% ============================================================

\subsection{Sphere recovers Mie}

For a homogeneous sphere of radius $a$ at $k_0 a\gg 1$ with refractive
index $\tilde n$, \cref{eq:RCS-PO} should reduce to the geometric-optics
limit of Mie theory: $\sigb(\kin,\kin)\to\pi a^2(1-T_0)$ for the
backward direction, $\sigb(\kin,-\kin)\to$ the forward-scatter peak
of width $\sim\lambda/a$.  The backscatter case is the trivial
specialisation of \cref{lem:peak} with $R_1=R_2=a$.  The forward
scatter is a Babinet/diffraction lobe and is computed in
\texttt{theory/scripts/mie\_theory\_corrected.py}; comparing the
PO/Kirchhoff prediction to the Mie reference at $a/\lambda\in[1,10]$
is a quick validation that the same residual decomposition of TAP
\S\,4 applies here.

\subsection{Cylinder, finite}

A finite cylinder of radius $a$ and length $L$ illuminated broadside
gives, again from \cref{eq:RCS-PO}, the textbook
$\sigb(\kin,\kin)=2\pi a L^2/\lambda\cdot(1-T_0)$ at the
specular ridge and $\sim a L\cdot(1-T_0)$ off-specular.  The
$L^2/\lambda$ enhancement is the body's geometric ``antenna gain''
along its axis.  For an upright torso ($a\approx\SI{15}{cm}$,
$L\approx\SI{60}{cm}$) at \SI{28}{GHz}, the broadside specular peak
is $\approx 1.0\,\mathrm{m^2}$ or $0\,$dBsm, matching the JSAC
paper's placeholder torso RCS.

\subsection{Mass body conservation}

A mass body of zero curvature would scatter nothing in the geometric
limit ($\kappa=0$ in \cref{eq:peak} gives infinite specular peak per
unit area times zero area times zero curvature; need to be careful
with the limit).  The right limit is a flat slab: PO predicts a single
specular reflection with reflectance $1-T_0$ and a Heaviside-shaped
beam.  The flat slab is the ``ground'' of any practical RT scene; the
RIB framework reduces to the standard slab-scattering treatment in
that limit.

% ============================================================
\section{Numerical sketch}\label{sec:numerical}
% ============================================================

The companion script
\texttt{theory/scripts/rib\_far\_field\_demo.py} carries out a small
numerical experiment.  Given the Thelonious phantom and the IT'IS
v5.0 skin tissue at \SI{28}{GHz}:

\begin{enumerate}[leftmargin=1.6em,itemsep=3pt]
\item Pick an incident plane wave with $\kin=-\hat z$
(downtilted, broadside) and unit polarisation
$\bm\psi=\hat x$.
\item Compute the per-triangle Fresnel reflection coefficients
$r_s,r_p$ on every front-facing triangle.
\item Evaluate the Kirchhoff integral \cref{eq:Fb} on a $64\times 64$
grid of outgoing directions $\kout$ on the upper hemisphere.
\item Plot $|\Frib(\kout)|^2$ on a Mollweide projection.
\item Verify the integrated identity
$\int_{S^2}|\Frib(\kout)|^2\diff\Omega/Z_0 \approx
\mathbf{x}^H\Qref\mathbf{x}$
where $\Qref$ is the reflection operator from
\texttt{q\_complement.tex} and $\mathbf{x}$ is the unit precoder
that places all power on the array element pointing at the body.
\end{enumerate}

The expected outcome: a strong specular lobe at $\kout=+\hat z$ of
the predicted angular width $\sim\SI{2.4}{\degree}$, a diffuse
background at $\sim 0.4$ of the specular peak, and a closed
energy budget within $\sim 5\%$ of the operator-level prediction.
The script prints both numbers.

If this works for a flat-incidence single plane wave, the same
machinery extends per-path to the full array: replace the single
$\kin$ with the multi-path superposition of \cref{eq:Erefl-array} and
the resulting bistatic pattern is the array's body-mediated radiation
pattern.

% ============================================================
\section{Loose threads and conjectures}\label{sec:threads}
% ============================================================

Same format as \texttt{q\_complement.tex}: questions, not claims.

\begin{enumerate}[leftmargin=1.6em,itemsep=3pt]
\item \emph{Closed-form bistatic RCS for a triaxial ellipsoid?}
The standard ellipsoid RCS formulae assume a perfect conductor and
have a closed form via the gradient of the surface phase.  The
Fresnel correction multiplies pointwise by $|r_0|^2=1-T_0$ (under
\ref{a:pb}), which is local and pulls outside the integral.  Hence:
$\sigb^{\mathrm{ellip}}(\kin,\kout)=
(1-T_0)\cdot\sigma_{\mathrm{PEC,ellip}}(\kin,\kout)$ should hold to
$5\%$.  An ellipsoid is a 3-parameter approximation of a torso.  This
gives a parametric RCS model with no integral.
\item \emph{Polarised bistatic.}  $\sigb$ above is unpolarised on the
outgoing side because $\Pi^{\mathrm{out}}$ is symmetric.  The
\emph{polarised} bistatic RCS $\sigb(\kin,\kout;\bm\psi^{\mathrm{in}}
,\bm\psi^{\mathrm{out}})$ has off-diagonal s/p terms that depend on
the local geometry of the specular point.  At pseudo-Brewster these
off-diagonals vanish; away from pseudo-Brewster they encode tissue
polarimetry.
\item \emph{Sensing-exposure beamspace duality.}  Form the
``backscatter beamspace'' as the eigenbasis of
$\overline{\Qbist}=\int\Qbist(\kout)\diff\Omega(\kout)=\Qref$ and
the ``exposure beamspace'' as the eigenbasis of $\Qabs$.  By
\texttt{q\_complement.tex} Prop~4.2 these are the same up to
rotation.  Conjecture: there is a finite-rank approximation
$\Qbist(\kout)\approx U\Lambda(\kout)U^H$ with $U$ the same shared
basis and $\Lambda(\kout)$ a diagonal angular-modulation matrix that
tells you ``which exposure mode lights up which receiver
direction''.  This is the body's beamspace transfer function; it
should be smooth in $\kout$ and possibly low-rank in mode index.
\item \emph{RIS analogy made literal.}  A uniform RIS of $N$
elements has $\Qbist^{\mathrm{RIS}}(\kout)=
\Phi^H\mathbf{a}(\kout)\mathbf{a}(\kout)^H\Phi$ with $\Phi$
diagonal-tunable.  A body has $\Qbist^{\mathrm{body}}(\kout)$ as in
\cref{def:Qbist} with $\Phi$ implicitly given by the surface phase
$e^{-ik_0\kout\cdot\rr}$ over the lit area.  The body is, in this
sense, a passive RIS with a fixed but pose-dependent phase pattern;
``re-pointing'' the body is achieved by re-posing it.  Quantifying
the angular-resolution this gives in the propagation graph is open.
\item \emph{Energy bookkeeping with the missed operator.}  The
direction-resolved book-keeping should read
$\int_{S^2}\Qbist(\kout)\diff\Omega+\Qabs+\Qmis=\mathbf{I}$, with
$\Qmis=\int_{S^2}\Qmis^{(\kout)}\diff\Omega$ the
direction-resolved missed operator from
\texttt{q\_complement.tex} Eq.~(15).  This is the natural array-level
generalisation of the optical theorem.  Cleaning up the proof is a
half-day job.
\item \emph{Pose differentiability.}  $\Frib^{(n)}(\kout)$ is a Kirchhoff
integral over $\Sigma$, which is parameterised by SMPL-X pose
parameters $\theta\in\Reals^{72}$.  $\partial\Frib^{(n)}/\partial\theta$
exists almost everywhere with the same GELU-smoothing trick as TAP
\S\,2.6.  This makes the body scattering pattern a
\emph{differentiable function of pose}, so RIB integrates with
gradient-based pose recovery from backscatter.  ``Imaging the body
from its own shadow'' becomes a tractable inverse problem.
\item \emph{Multi-body RIB capacity bound.}  In a plaza with $K$
bystanders, what is the channel-rank lift $\Delta r$ from the
RIB-network reflectors?  Conjecture: $\Delta r\sim K\cdot
\mathrm{rank}\overline{\Qbist}/M$ in the dense-crowd limit; for
\SI{28}{GHz} and a $K=50$ plaza this is $\sim 50\cdot 5/64\sim 4$
extra ranks per receiver, which shifts the JSAC paper's observed
effective-rank-$4$ result from a curiosity into a derivable number.
\item \emph{Empirical RCS calibration.}  $\sigb$ at the body is a
quantity that mmWave radar literature has measured, dimensionally,
many times.  Cross-checking the PO prediction \cref{eq:RCS-PO}
against e.g.\ Geng \emph{et al.} 2018 ``Human Body Radar Cross
Section'' (28\,GHz, single subject, $\sim$0\,dBsm at broadside)
should give a one-figure validation.
\end{enumerate}

% ============================================================
\section{Three observables, one operator family}\label{sec:triple}
% ============================================================

A consequence of the construction worth pulling out separately, even
if it is a closing remark rather than a load-bearing theorem, is
that the same operator family
$\{\Qabs,\Qref,\Qbist(\kout)\}$ governs the body's three
electromagnetic observables seen from the array:

\begin{tabular}{ll}
\textbf{Transmit, absorbed:} &
$P_{\mathrm{ab}}(\mathbf{x})=\mathbf{x}^H\Qabs\mathbf{x}$ \\
\textbf{Transmit, scattered to direction $\kout$:} &
$P_{\mathrm{re}}(\mathbf{x};\kout)=\mathbf{x}^H\Qbist(\kout)\mathbf{x}$ \\
\textbf{Receive, body thermal pickup:} &
$\langle|\mathbf{y}^H\mathbf{n}_{\mathrm{body}}|^2\rangle
\propto\mathbf{y}^H\Qabs\mathbf{y}$ \\
\end{tabular}

\noindent The first two are the two sides of the
\texttt{q\_complement.tex} energy ledger plus its angular
refinement.  The third is the Kirchhoff dual of
\texttt{q\_complement.tex} Prop.~6.1.  Two consequences of having
all three controlled by the same shared eigenbasis (per
\texttt{q\_complement.tex} Prop~4.2 and the per-direction extension
\cref{prop:perdir}):

\begin{enumerate}[leftmargin=1.6em,itemsep=2pt]
\item \textbf{Cross-validation.}  The DTN's prediction of any one
of the three observables is testable against any other.  Predict
$\Qabs$ from pose; verify by listening to body thermal radiation on
idle radios; cross-check by measuring backscatter under a chosen
precoder.  Three different physical channels, one model.
\item \textbf{Phantom-free compliance.}  The trace
$\tr\Qabs$ is recoverable from the per-element noise
PSD of the array under thermal load (passive measurement, no
phantom needed); the largest eigenvalue $\lambda_1(\Qabs)$ requires
an interferometric receive-side measurement (cross-element noise
correlations).  Both can be done in situ on a deployed array, in
principle.  Whether this is practical depends on the noise floor
relative to body thermal pickup at $\sim$\SI{300}{K}, which at
\SI{28}{GHz} and a body filling the array's main beam gives
$\sim$\SI{-30}{dBm} pickup against a typical $\sim$\SI{-90}{dBm}
thermal noise floor; comfortably above noise.  This is a regulatory
audit story I do not yet know if regulators would entertain, but it
is technically defensible as far as I can see.
\end{enumerate}

The point is conceptual: a body in front of a JSAC array is one
electromagnetic object, not three.  The DTN that tracks one of its
observables can predict the others.

% ============================================================
\section{Summary}\label{sec:sum}
% ============================================================

The pitch, in three sentences.  The JSAC body twin already builds
$\Qabs^{(u)}$ via a surface integral over the body mesh; one
extra direction-resolved phase factor on the same loop gives
$\Qbist^{(u)}(\kout)$, the body's bistatic scattering matrix, and
$\sigb^{(u)}(\kin,\kout)$, its analytical bistatic RCS.  The
per-direction pseudo-Brewster locking
$\Qbist(\kout)\approx\frac{1-T_0}{T_0}\Qabs^{(\kout)}$ is what
makes the rest of the DT story interesting: capacity gain from the
RIB-network is hard-bounded by the array's exposure dose, the
empirical effective rank $3{-}4$ has a model-based explanation, and
calibration follows the same ridge-LS recipe the JSAC paper already
uses for path amplitudes.

In the propagation graph a body, treated this way, is not a blocker
and not an opaque box.  It is a self-tuning passive RIS whose only
knob is gait, whose pattern is computable from pose, mesh, tissue,
and frequency, and whose three observables (absorbed, scattered,
thermal) factor through one shared operator family.  The
intelligence is the DTN's; the body just stands there.  The cost of
RIB to a DTN that is already running an exposure twin is one phase
factor per outgoing direction per body per slot, which is to say
nothing.  The benefit is everything the previous sections argued for.

What this is not, and where the story lives next.  RIB is a
mmWave-and-above tool; below \SI{6}{GHz} the body is a wavelength-scale
lossy blob, PO does not apply, and the right framework is full-wave.
The intended home of these ideas is the JSAC SI submission on
Digital Twins for Wireless Networks; specifically, \cref{sec:dtn} and
\cref{thm:pareto} are the additions that should land in the JSAC
paper, while the remainder (Kirchhoff dual, transparency subspace,
pose-differentiable RCS) is more naturally a follow-up note.

\end{document}
